\documentclass[11pt,a4paper]{article}
\usepackage[margin=25mm]{geometry}
\usepackage{amsmath,amssymb,amsthm}
\usepackage{hyperref}
\hypersetup{colorlinks=true,urlcolor=blue,linkcolor=blue}
\newtheorem{theorem}{Theorem}
\newtheorem{lemma}{Lemma}
\newcommand{\diag}{\operatorname{diag}}
\newcommand{\rad}{\rho}
\title{Generic local identifiability does not imply generic global identifiability\texorpdfstring{\\}{ }even when $\rho(G)<1$}
\author{Research note prepared with Codex}
\date{8 October 2026}
\begin{document}
\maketitle
\begin{abstract}
We give a twelve-node network with zero diagonal and forty-two independently
unknown edge entries. Nine nodes are excited and nine are measured by fixed
selection matrices. The observation map $G\mapsto C(I-G)^{-1}B$ is generically
locally identifiable. A nonempty open subset of the admissible parameter
space has two distinct compatible networks, both satisfying $\rho(G)<1$.
The counterexample therefore disproves generic local implying generic global
identifiability even with this spectral-radius restriction. An explicit
rational pair and exact matrix-power bounds certify the claim. The model
uses no prescribed equalities between unknown edges, no known nonzero
edges, and no self-loops.
\end{abstract}

\paragraph{Scope and the meaning of generic.}
We consider $w=Gw+Br$, $y=Cw$, with known binary node selections $B,C$,
zero diagonal of $G$, and independently unknown entries at its permitted
off-diagonal positions. Impose $\rho(G)<1$ on every admissible candidate,
including any alternative network. The stronger assumption appears in
the later Legat--Hendrickx manuscript \cite{LH24}.
The matrices in the explicit pair are real constants, so they also give
proper, stable rational transfer functions satisfying the spectral-radius
restriction at every frequency. In the supplied paper's notation,
$P=I-G$, $Q=B$, $R=C$, $S=0$.

Generic local identifiability means that the observation Jacobian has full
column rank outside a proper algebraic subset. Generic global identifiability
on the restricted domain would require uniqueness for almost every
admissible parameter. A nonempty Euclidean open set of ambiguous admissible
parameters disproves this property. We do not claim that every stable
parameter in the structure has two stable alternatives.

\paragraph{The six-node ambiguity used in the construction.}
Let
\[
 A=\begin{pmatrix}0&0&a_3\\a_1&0&0\\0&a_2&0\end{pmatrix},\quad
 D=\begin{pmatrix}0&d_1&0\\0&0&d_2\\d_3&0&0\end{pmatrix},\quad
 X=\diag(x_1,x_2,x_3),
\]
and
\[
 G_b=\begin{pmatrix}A&0\\X&D\end{pmatrix},\qquad
 B_b=\begin{pmatrix}I_3\\0\end{pmatrix},\qquad
 C_b=\begin{pmatrix}0&I_3\end{pmatrix}.
\]
The nine quantities $a_i,d_i,x_i$ are independent. Its observation is
\[
 F=C_b(I_6-G_b)^{-1}B_b=(I_3-D)^{-1}X(I_3-A)^{-1}.
\]
Set $u_i=a_i/x_i$, $v_i=d_i/x_i$, $z_i=1/x_i$.
The six off-diagonal entries of $F^{-1}$ determine $u_i,v_i$ directly,
and its diagonal entries are
\begin{equation}\label{eq:diags}
 e_i=z_i+\frac{u_{i-1}v_{i-1}}{z_{i-1}},\qquad i\pmod 3.
\end{equation}
With $u,v$ fixed, the determinant of the Jacobian of these three equations
with respect to $z$ is
\begin{equation}\label{eq:local}
 \det J_z=1-\prod_{i=1}^3 a_id_i.
\end{equation}
It is nonzero generically. The three equations in \eqref{eq:diags} compose
to a M\"obius fixed-point equation, hence a quadratic for $z_1$.

For completeness, a rational formula for the alternative is
\begin{equation}\label{eq:branch}
 k_i=a_id_i\frac{x_{i+1}}{x_i},\quad
 h_i=1+k_i+k_ik_{i+1},\quad
 \eta_i=\frac{k_ih_{i+1}}{h_i},\quad
 (a_i',d_i',x_i')=\frac{(a_i,d_i,x_i)}{\eta_i}.
\end{equation}
The transformation leaves $u_i,v_i$ fixed and sends $z_i$ to $\eta_i z_i$.
Writing $L=\prod_i k_i=\prod_i a_id_i$, the identities
\[
 \eta_i-1=\frac{L-1}{h_i},\qquad
 \eta_i+\frac{k_{i-1}}{\eta_{i-1}}=1+k_{i-1}
\]
show that it preserves all observations and gives a distinct second point
when $L\ne1$. It is well-defined and continuous near the points used below.

\paragraph{Augmentation that preserves all fibers.}
Add six nodes with signal $h$, all of which are excited and measured.
Allow feedback through matrices $K\in\mathbb R^{3\times6}$,
$L_c\in\mathbb R^{6\times3}$ and $J\in\mathbb R^{6\times6}$:
\begin{equation}\label{eq:aug}
 \mathcal G=\begin{pmatrix}G_b&B_bK\\L_cC_b&J\end{pmatrix},\qquad
 \mathcal B=\begin{pmatrix}B_b&0\\0&I_6\end{pmatrix},\qquad
 \mathcal C=\begin{pmatrix}C_b&0\\0&I_6\end{pmatrix}.
\end{equation}
Thus $\mathcal B$ excites nodes $1,2,3,7,8,9,10,11,12$, and
$\mathcal C$ measures nodes $4,5,6,7,8,9,10,11,12$.

\begin{lemma}\label{lem:inverse}
Where the displayed inverses exist, the augmented observation
$H=\mathcal C(I_{12}-\mathcal G)^{-1}\mathcal B$ satisfies
\begin{equation}\label{eq:inverse}
 H^{-1}=\begin{pmatrix}F^{-1}&-K\\-L_c&I_6-J\end{pmatrix}.
\end{equation}
Consequently, the observation determines every permitted entry of
$K,L_c,J$, and its remaining ambiguity is precisely the ambiguity in $G_b$.
\end{lemma}
\begin{proof}
Let the original and new excitations be $r$ and $s$. The node equations imply
$y=F(r+Kh)$ and $h=L_cy+Jh+s$. Rearranging gives
\[
 \begin{pmatrix}r\\s\end{pmatrix}
 =\begin{pmatrix}F^{-1}&-K\\-L_c&I_6-J\end{pmatrix}
  \begin{pmatrix}y\\h\end{pmatrix},
\]
which proves the identity and its identifiability consequences.
\end{proof}

\paragraph{An explicit pair with small spectral radii.}
Let
\[
 U=\begin{pmatrix}0&0&1\\1&0&0\\0&1&0\end{pmatrix},\qquad
 V=U+U^T,\qquad I=I_3.
\]
Choose the numerical values
\[
 K_0=\begin{pmatrix}V&-U\end{pmatrix},\qquad
 L_0=\begin{pmatrix}U^T\\-V\end{pmatrix},\qquad
 J_0=\begin{pmatrix}U^T&-V\\I&U\end{pmatrix}.
\]
For scalars $r,x$, define the following matrix in $3\times3$ blocks:
\begin{equation}\label{eq:explicit}
 M(r,x)=\begin{pmatrix}
 -rU&0&V&-U\\
 xI&-rU^T&0&0\\
 0&U^T&U^T&-V\\
 0&-V&I&U
 \end{pmatrix}.
\end{equation}
Now take
\begin{equation}\label{eq:pair}
 r=\frac{99}{100},\qquad
 \mathcal G_1=M(r,1),\qquad
 \mathcal G_2=M(r^{-1},r^{-2}).
\end{equation}
Both matrices have exactly the same zero pattern and zero diagonal. The
model assigns a separate independent unknown to \emph{each} of the forty-two
permitted entries of this pattern: nine in $G_b$, nine in $K$, nine in $L_c$
and fifteen in $J$. The block notation in \eqref{eq:explicit} gives convenient
values for two points of this full model; it imposes no equality constraints
on the model's unknowns and fixes no nonzero edge in advance.

\begin{theorem}\label{thm:main}
The forty-two-parameter network structure just defined is generically
locally identifiable but is not generically globally identifiable on the
domain $\rho(\mathcal G)<1$.
\end{theorem}
\begin{proof}
\emph{Equality of observations.}
For the original six-node block in $M(r,x)$,
\[
 F(r,x)^{-1}=x^{-1}(I+rU)(I+rU^T)
 =x^{-1}\big[(1+r^2)I+rV\big].
\]
Therefore $F(r,1)=F(r^{-1},r^{-2})$. Since both augmented networks use
the same $K_0,L_0,J_0$, Lemma~\ref{lem:inverse} proves that their complete
$9\times9$ observations are equal. In fact their common inverse observation is
\[
 H^{-1}=\begin{pmatrix}
 (1+r^2)I+rV&-V&U\\
 -U^T&I-U^T&V\\
 V&-I&I-U
 \end{pmatrix}.
\]

\emph{Exact spectral-radius certificates.}
Direct rational multiplication of the explicit matrices gives
\begin{equation}\label{eq:norms}
 \|\mathcal G_1^4\|_\infty=\frac{33910001}{100000000}<1,
 \qquad
 \|\mathcal G_2^4\|_\infty=\frac{16976126}{32019867}<1.
\end{equation}
Here $\|\cdot\|_\infty$ is the maximum absolute row sum. Hence
\[
 \rho(\mathcal G_i)^4=\rho(\mathcal G_i^4)
 \le\|\mathcal G_i^4\|_\infty<1,
\]
proving the required restriction for both candidates without relying on
floating-point eigenvalues. Their numerical spectral radii are approximately
$0.509563$ and $0.404210$, respectively.

\emph{Generic local identifiability.}
In \eqref{eq:inverse}, the thirty-three controller entries occur separately
and linearly. The remaining $3\times3$ block $F^{-1}$ has generic Jacobian
rank nine by \eqref{eq:local}. Consequently the full observation Jacobian
has generic rank $9+33=42$. It has rank forty-two at both explicit points,
since $\prod_i a_id_i=r^6$ or $r^{-6}$, neither of which is one.

\emph{The ambiguity occurs on an open admissible set.}
Apply \eqref{eq:branch} to the nine independent entries of $G_b$, and keep
the thirty-three independent entries of $K,L_c,J$ unchanged. At the first
point in \eqref{eq:pair}, this maps to the second point. The formulas are
rational and continuous on a neighborhood of that point. The two networks
are distinct and both inequalities \eqref{eq:norms} are strict, so a
sufficiently small open neighborhood in the full forty-two-dimensional
parameter space consists of distinct pairs satisfying $\rho<1$ and having
identical observations. This set has positive measure and cannot be
contained in a proper algebraic subset. Thus generic global identifiability
fails even after restricting both candidates to $\rho<1$.
\end{proof}

\paragraph{Why the augmentation can make both branches stable.}
The two six-node branches coalesce at $r=1$. Their common augmented matrix
satisfies
\begin{equation}\label{eq:nilpotent}
 M(1,1)^4=0.
\end{equation}
This can be understood without multiplying twelve-dimensional matrices.
At $r=x=1$, let $G_{b0}$ denote the six-node block. The controller satisfies
$J_0=G_{b0}+B_bK_0-L_0C_b$. A change of coordinates from $(w,h)$ to
$(w,w-h)$ makes the augmented matrix block upper triangular, with diagonal
blocks
\[
 G_{b0}+B_bK_0=\begin{pmatrix}U^T&-U\\I&-U^T\end{pmatrix},\qquad
 G_{b0}-L_0C_b=\begin{pmatrix}-U&-U^T\\I&U\end{pmatrix}.
\]
Each of these squares to zero, proving \eqref{eq:nilpotent}. By continuity,
both $M(r,1)$ and $M(r^{-1},r^{-2})$ have arbitrarily small spectral radii
when $r\ne1$ is sufficiently close to one. The ambiguity can therefore
occur with $\rho<\varepsilon$ for any prescribed $\varepsilon>0$.
At $r=1$ the local rank drops; the counterexamples use $r\ne1$.

\paragraph{Reproducible checks.}
The script \texttt{verify\_spectral\_counterexample.py} checks:
the common zero-diagonal pattern and edge count;
equality of the complete rational $9\times9$ observations;
the block inverse identity;
the exact norms in \eqref{eq:norms}; and nilpotence in \eqref{eq:nilpotent}.
It also forms the full $81\times42$ observation Jacobian and certifies rank
forty-two at both points by elimination modulo the prime $1000003$.
A nonzero minor modulo this prime certifies a nonzero rational minor.
All denominators used in this certificate are invertible modulo the prime.

As a check against accidental equality constraints, the script independently
perturbs all forty-two edge values to distinct rational numbers, applies
the rational alternative formula, and verifies exact observation equality
and both strict matrix-power bounds again. The detailed results and explicit
matrices are saved in \texttt{spectral\_verification\_results.json}.
The proof above is independent of the numerical search used while exploring
other possible constructions. No minimality or literature-priority claim
is made.

\begin{thebibliography}{9}
\bibitem{LH20}
A. Legat and J. M. Hendrickx,
``Local Network Identifiability with Partial Excitation and Measurement,''
\emph{59th IEEE Conference on Decision and Control}, 2020, pp. 4342--4347.
\href{https://perso.uclouvain.be/julien.hendrickx/availablepublications/LH2020_CDC.pdf}{Author's PDF}.

\bibitem{LH24}
A. Legat and J. M. Hendrickx,
``Identifiability in Networked Systems: Algebraic and Combinatorial
Conditions,'' manuscript listed by the author as submitted in August 2024.
\href{https://perso.uclouvain.be/antoine.legat/_static/TAC24.pdf}{Author's PDF}.
\end{thebibliography}
\end{document}
